\section{Curvature, diffusion, and entropy calculus}\label{sec:calculus}

We first establish the analytic facts used to construct and study the
probability measures in the proof. Two time conventions will be kept
distinct: convolution with a Gaussian of variance $r$ has generator
$\Delta/2$, whereas a stationary diffusion has generator $-A$ and noise
coefficient $\sqrt2$. These conventions determine every entropy
dissipation factor below.

For probabilities $p,\pi$ with $p\ll\pi$, define their relative entropy by
\[
 H(p\mid\pi)=\int\log\frac{dp}{d\pi}\,dp.
\]
Set $H(p\mid\pi)=+\infty$ when absolute continuity fails. Write
$\gamma$ for standard Gaussian measure in the relevant dimension, and set
\[
 \Ical(p\mid\pi)=4\int\left|D\sqrt{\frac{dp}{d\pi}}\right|^2d\pi.
\]
The value is $+\infty$ unless the density exists and its square root
belongs to the weighted Sobolev space. When the densities are positive
and smooth this is $\int|D\log(dp/d\pi)|^2dp$. Our convention is
\[
 \Ent_\pi(f^2)\le 2R(\pi)\int|Df|^2d\pi,
 \qquad 2H(p\mid\pi)\le R(\pi)\Ical(p\mid\pi).
\]
The Poincar\'e constant $R^P(\pi)$ uses the convention
\[
 \Var_\pi(f)=\int f^2d\pi-\left(\int f\,d\pi\right)^2
 \le R^P(\pi)\int|Df|^2d\pi.
\]
Linearization of the entropy inequality gives $R^P(\pi)\le R(\pi)$.

For smooth positive densities the Sobolev space is the completion of
$C_c^\infty$ in the squared norm $\int(f^2+|Df|^2)d\pi$.
Cutoffs, truncation of the values, and smoothing on compact sets show
that it also consists of the functions with a weak gradient in this
norm. The same approximation and lower semicontinuity of entropy
extend a compact-test logarithmic Sobolev inequality to this space.
One convenient proof of the latter lower semicontinuity is the formula
\begin{equation}\label{eq:calculus-entropy-duality}
 H(p\mid\pi)=\sup_{g\ \mathrm{bounded}}
 \left\{\int g\,dp-\log\int e^g d\pi\right\}.
\end{equation}
Jensen's inequality proves the upper bound; bounded truncations of
$\log(dp/d\pi)$ give equality. The same formula implies data
processing under a measurable map or a Markov kernel, by conditional
Jensen. Conditional entropy identities below are always averaged over
the conditioning variable. In particular, disintegrating a joint
density and splitting its logarithm gives the chain rule
\[
 H(P_{U,Z}\mid Q_U\otimes\nu)
 =H(P_U\mid Q_U)+\E_{P_U}H(P_{Z\mid U}\mid\nu).
\]
It holds in the extended nonnegative sense by truncation. Mutual
information is $\MI(U;Z)=H(P_{U,Z}\mid P_U\otimes P_Z)$.
Taking $g(z)=v\cdot z$ in Equation~\eqref{eq:calculus-entropy-duality}
and optimizing over $v$ also gives
$H(p\mid\gamma)\ge|\E_pZ|^2/2$ when the mean exists.

The following is the anisotropic Bakry--\'Emery curvature criterion
\cite{BakryEmery1985}.  We include its semigroup proof and the convex
approximation needed for Gaussian posteriors with nonsmooth support.

\begin{lemma}[Positive curvature]\label{lem:curvature}
Let $S$ be a positive definite matrix on $\R^n$, and let $U$ be a
proper lower semicontinuous convex function, possibly taking the value
$+\infty$. Suppose
\[
 d\pi(x)=Z^{-1}\exp\{-\tfrac12\langle x,Sx\rangle-U(x)\}\,dx
\]
is a probability with a full dimensional support. Then
\[
 \Ent_\pi(f^2)\le2\E_\pi\langle Df,S^{-1}Df\rangle,
 \qquad
 \Var_\pi(f)\le\E_\pi\langle Df,S^{-1}Df\rangle.
\]
The inequalities hold first for compact smooth tests, and then for
functions obtained by finite-energy approximation.
\end{lemma}

\begin{proof}
A linear change of variables reduces the assertion to $S=\Id$.
First suppose $V(x)=|x|^2/2+U(x)$ is smooth with bounded Hessian.
The diffusion with generator $L=\Delta-DV\cdot D$ has globally
Lipschitz drift and invariant law $\pi$. If $J_t$ is the derivative
of its trajectory with respect to the starting point, then
$\dot J_t=-D^2V(X_t)J_t$, whence $\|J_t\|_{\op}\le e^{-t}$.
For smooth bounded positive $q$, bounded away from zero,
weighted Cauchy--Schwarz gives
\[
 \frac{|DP_tq|^2}{P_tq}
 \le e^{-2t}P_t\left(\frac{|Dq|^2}{q}\right).
\]
Integration by parts yields
\[
 \frac{d}{dt}\int (P_tq)\log(P_tq)\,d\pi
 =-\int\frac{|DP_tq|^2}{P_tq}\,d\pi.
\]
Couple a trajectory started at $x$ to an independent stationary
initial point with the same driving Brownian motion. Their distance
is at most $e^{-t}$ times their initial distance. Since $\pi$ has
Gaussian tails, this proves convergence of $P_tq$ to $\E_\pi q$
for bounded Lipschitz $q$, and hence convergence of its entropy to
zero. Integrating the preceding two displays gives
$\Ent_\pi(q)\le\frac12\int|Dq|^2/q\,d\pi$. Apply this to
$q=f^2+\varepsilon$ and let $\varepsilon\downarrow0$.

For general $U$, take its quadratic infimal convolutions
\[
 U_j(x)=\inf_z\{U(z)+j|x-z|^2\}.
\]
The minimizer is unique, and monotonicity of the subgradient of $U$
shows that $DU_j$ is Lipschitz. The functions $U_j$ are convex and
increase to $U$. A fixed affine minorant $U(z)\ge\ell\cdot z+c$
gives $U_j(x)\ge\ell\cdot x+c-|\ell|^2/(4j)$.
Mollify $U_j$ with a smooth compactly supported kernel at radii
tending to zero. The resulting functions are smooth, convex, and
have bounded Hessian; the affine lower bound persists with a
uniform additive error. Consequently their normalized densities
converge almost everywhere and in $L^1$, dominated by an integrable
Gaussian times an exponential of a linear function. The compact-test
inequality passes to the limit. Boundary values of $U$ do not matter,
since the boundary of a full dimensional convex support has
Lebesgue measure zero. Finally, insert $f=1+\varepsilon g$,
using finite-energy cutoffs of the constant, and compare the
coefficients of $\varepsilon^2$ to obtain the variance inequality.
\end{proof}

We shall also use two immediate consequences with their constants.
If independent laws satisfy inequalities with quadratic energy
forms $Q_1,Q_2$, their product satisfies the inequality with energy
$Q_1+Q_2$. Indeed, decompose entropy into conditional entropy and
the entropy of the fiber integral, and use
\[
 \left|D_x\left(\int f(x,y)^2d\pi_2(y)\right)^{1/2}\right|^2
 \le\int|D_xf(x,y)|^2d\pi_2(y).
\]
Approximation removes zeros of the fiber norm. The variance version
uses conditional variance and differentiation of the fiber mean.
In particular, independent addition gives
\begin{equation}\label{eq:calculus-convolution}
 R(\Law(X+Z))\le R(\Law(X))+R(\Law(Z)).
\end{equation}
Second, Herbst's argument (see \cite[Section~2.3]{Ledoux1999})
shows that, if $R(\pi)\le R$, every integrable $1$-Lipschitz function
$u$ satisfies
\begin{equation}\label{eq:calculus-herbst}
 \log\E_\pi e^{t(u-\E_\pi u)}\le Rt^2/2.
\end{equation}
For a bounded smooth $u$, applying the inequality to $e^{tu/2}$,
and writing $\psi(t)=\log\E e^{tu}$, gives
$t\psi'(t)-\psi(t)\le Rt^2/2$. Integrate the derivative of
$\psi(t)/t$ from zero. Bounded Lipschitz truncations and Fatou's
lemma prove the assertion in general. In particular the conclusion
applies to linear forms whenever their first moments are finite.

\begin{lemma}[Dirichlet operator and stationary diffusion]
\label{lem:dirichlet-core}
Let $\pi$ be a smooth, strictly positive probability density on
$\R^n$. On $C_c^\infty(\R^n)$ define
\[
 A_\pi=-\Delta-D\log\pi\cdot D.
\]
This nonnegative symmetric operator is essentially self-adjoint in
$L^2(\pi)$, and $C_c^\infty$ is a core for its closure. In
particular, if $f\in L^2(\pi)$ and the distribution $A_\pi f$
belongs to $L^2(\pi)$, then $f\in\Dom(A_\pi)$. Its semigroup
$P_u=e^{-uA_\pi}$ preserves constants and is realized by the
stationary diffusion
\begin{equation}\label{eq:calculus-diffusion}
 dY_u=\sqrt2\,dB_u+D\log\pi(Y_u)\,du,
 \qquad Y_0\sim\pi.
\end{equation}
For $f\in\Dom(A_\pi)$ the stationary It\^o formula holds in $L^2$,
with drift $-A_\pi f$ and martingale integrand $\sqrt2\,Df$.
\end{lemma}

\begin{proof}
Integration by parts gives
$\langle f,A_\pi g\rangle=\int Df\cdot Dg\,d\pi$.
If $v\in L^2(\pi)$ is orthogonal to
$(1+A_\pi)C_c^\infty$, then $(1+A_\pi)v=0$
distributionally.  Conjugating by $\sqrt\pi$ reduces this to a
Poisson equation with a smooth local potential.  The local Newtonian
potential estimates and bootstrap make $v$ smooth; see
\cite[Lecture~24, p.~2; Lecture~16, p.~2]{Viaclovsky2004}. Take a compact cutoff $\chi_R$, equal to one on the ball
of radius $R$, with $|D\chi_R|\le C/R$. Testing by
$\chi_R^2v$ and using $2ab\le a^2/2+2b^2$ gives
\[
 \int\chi_R^2v^2d\pi+\tfrac12\int\chi_R^2|Dv|^2d\pi
 \le2\int|D\chi_R|^2v^2d\pi.
\]
Thus $v=0$. The range of $1+A_\pi$ is dense. Nonnegativity
gives $\|(1+A_\pi)f\|_2\ge\|f\|_2$, so the range of the
closure is also closed and therefore equals $L^2$. The closure is
self-adjoint: for $g$ in its adjoint domain, solve
$(1+\overline A)h=(1+A^*)g$; then $g-h$ is in the orthogonal
complement of the range of $1+A$, hence vanishes. This also proves
the asserted maximal distributional domain characterization.

The cutoffs $\chi_R$ converge to $1$ in the form norm with energy
tending to zero. Hence $1$ belongs to the nullspace of the closed
form and $P_u1=1$. To identify the diffusion, solve the stochastic
equation locally and kill it on exiting each ball. Its kernels are
symmetric and sub-Markov, as follows either from the Dirichlet
heat equation or integration by parts. Increasing the balls gives
the minimal sub-Markov semigroup. The stopped It\^o formula on
compact tests identifies its generator as an extension of $-A_\pi$.
Essential self-adjointness identifies it with $P_u$, which loses
no mass. Started in $\pi$, the resulting diffusion is stationary
and nonexplosive. Finally approximate $f$ in graph norm by compact
smooth functions. Their gradients converge in $L^2(\pi)$ since
$\|Dg\|_2^2=\langle g,Ag\rangle$. Stationarity and the
It\^o isometry pass every term of their formulas to the limit.
\end{proof}

\begin{lemma}[Entropy dissipation]\label{lem:entropy-calculus}
For the semigroup of Lemma~\ref{lem:dirichlet-core}, if
$p_0=q_0\pi$ has finite relative entropy and
$p_u=(P_uq_0)\pi$, then, for $0\le a\le b$,
\begin{equation}\label{eq:calculus-stationary-entropy}
 H(p_a\mid\pi)-H(p_b\mid\pi)
 \ge\int_a^b\Ical(p_u\mid\pi)\,du.
\end{equation}
If instead both probability measures evolve by Gaussian convolution,
$p_r=p_0*\gamma_r$ and $q_r=q_0*\gamma_r$, then
\begin{equation}\label{eq:calculus-heat-entropy}
 H(p_a\mid q_a)-H(p_b\mid q_b)
 \ge\frac12\int_a^b\Ical(p_r\mid q_r)\,dr
\end{equation}
whenever the entropy on the left at time $a$ is finite. Here
$\gamma_r$ is the centered Gaussian with covariance $r\Id$,
and time zero is allowed. For the standard Gaussian
Ornstein--Uhlenbeck semigroup, Equation~\eqref{eq:calculus-stationary-entropy}
is an equality, including time zero when the initial entropy is
finite. Each assertion also holds after integration over conditioning
labels whose noises are independent conditional on those labels.
\end{lemma}

\begin{proof}
For a bounded density $q$ bounded away from zero, the form chain
rule, first for smooth functions and then by form approximation,
gives
\[
 \frac{d}{du}\int(P_uq)\log(P_uq)\,d\pi
 =-\int\frac{|DP_uq|^2}{P_uq}\,d\pi.
\]
Spectral smoothing puts $P_uq$ in the operator domain for $u>0$;
the bounded range permits $1+\log(P_uq)$ as a form test.
Approximate an arbitrary finite-entropy density by truncating it
above and below and normalizing. The initial entropies converge.
The semigroup contracts $L^1$, so densities converge at every
later time. Lower semicontinuity of entropy and of
$4\int|D\sqrt q|^2d\pi$, followed by Fatou's lemma in time,
gives Equation~\eqref{eq:calculus-stationary-entropy} from time zero.
Restarting the argument gives arbitrary $a$.

For simultaneous heat evolution put $z=p_r/q_r$. Direct
differentiation, using $\partial_r p_r=\Delta p_r/2$ and the
same equation for $q_r$, gives
\begin{equation}\label{eq:calculus-heat-chain}
 (\partial_r-\tfrac12\Delta)\{q_r\Phi(z)\}
 =-\tfrac12q_r\Phi''(z)|Dz|^2.
\end{equation}
Here use smooth convex $\Phi_k$ with
$\Phi_k(1)=\Phi_k'(1)=0$, with second derivatives increasing to
$1/z$, and with second derivatives supported in compact subsets
of $(0,\infty)$. Such functions have at most linear growth and
increase to $z\log z-z+1$. Integrate
Equation~\eqref{eq:calculus-heat-chain} against $\chi_R$ with
$|\Delta\chi_R|\le C/R^2$. At each fixed $k$ the cutoff error
vanishes because $q_r\Phi_k(p_r/q_r)\le C_k(p_r+q_r)$.
Let $R\to\infty$ and then $k\to\infty$. Monotone convergence
gives the entropy and information terms, and in particular the
stated inequality on positive time intervals. Data processing
gives $H(p_\varepsilon\mid q_\varepsilon)\le H(p_0\mid q_0)$;
let $\varepsilon\downarrow0$ to include the initial time.

For Ornstein--Uhlenbeck evolution relative to $\gamma$, apply
the same convex truncations to the ratio. Its generator is
$\Delta-x\cdot D$. Expanding radial cutoffs have bounded
generator, tending pointwise to zero, so the error vanishes by
domination at fixed $k$. This gives equality on positive time
intervals. At time zero, relative entropy is continuous: data
processing gives its upper bound and lower semicontinuity under
weak convergence gives its lower bound. Monotone convergence of
the nonnegative time integral includes zero. Tonelli's theorem
proves all conditional versions.
\end{proof}

We use the relative-entropy form of the Gaussian-channel identity of
Guo, Shamai, and Verd\'u \cite[Theorem~5, Equation~(66)]{GuoShamaiVerdu2005}.
The proof records explicitly its conditional integrated form.

\begin{lemma}[Gaussian signals]\label{lem:gaussian-signals}
Let $S$ be an $\R^d$-valued random vector with finite second
moment, let $U$ be any side information, and let $G$ be standard
Gaussian independent of $(S,U)$. Write $Z_t=\sqrt tS+G$ and
\[
 d(t)=\E_U H(\Law(Z_t\mid U)\mid\gamma).
\]
Then $d(0)=0$, $d(t)\le t\E|S|^2/2$, and
\begin{equation}\label{eq:calculus-gaussian-signal}
 d(t)=\frac12\int_0^t
 \E\left|\E[S\mid Z_v,U]\right|^2dv.
\end{equation}
The integrand is nondecreasing in $v$. The formula without side
information is obtained by taking $U$ constant.
\end{lemma}

\begin{proof}
Conditional convexity of relative entropy and
$H(\gamma(\cdot-\sqrt t s)\mid\gamma)=t|s|^2/2$
give the upper bound. Gaussian kernel differentiation gives
\[
 D_z\log\frac{d\Law(Z_t\mid U)}{d\gamma}(z)
 =\sqrt t\,\E[S\mid Z_t=z,U].
\]
The differentiation is justified at each positive $t$ by the
Gaussian kernel and the finite second moment, or by truncating
$S$ and passing locally and in the conditional squared norm.
Applying Ornstein--Uhlenbeck evolution for time $u$ replaces $t$
by $te^{-2u}$. Its entropy dissipation identity and
$dt/du=-2t$ give $d'(t)=\frac12\E|\E[S\mid Z_t,U]|^2$
almost everywhere on $(0,\infty)$. The upper bound implies
$d(t)\to0$ at zero, proving the integrated identity. For
$0<v<t$, the channel at $v$ is obtained from the channel at $t$
by multiplying its output by $\sqrt{v/t}$ and adding fresh
Gaussian noise of covariance $(1-v/t)\Id$. Conditional Jensen
then proves monotonicity of prediction energy.
\end{proof}

\subsection{The transport--entropy consequence}
\label{sec:transport-consequence}

\begin{proof}[Proof of Corollary~\ref{cor:intro-transport}]
Jensen's inequality in each coordinate gives
$\int|x|^2d\mu\le na^2\log2$, so $\mu\in\mathcal P_2(\R^n)$.
Only the case $H(\nu\mid\mu)<\infty$ needs proof. For $s>0$, put
$\gamma_s=N(0,s\Id)$, $\mu_s=\mu*\gamma_s$, and
$\nu_s=\nu*\gamma_s$. The density of $\mu_s$ is smooth and strictly
positive. For $X\sim\mu$ and an independent standard Gaussian $G$,
its potential $V_s$ satisfies
\[
 D^2V_s(y)=s^{-1}\Id-s^{-2}\Cov(X\mid X+\sqrt sG=y)\ge0.
\]
Indeed, the posterior has the original convex potential plus
$|x-y|^2/(2s)$, so Lemma~\ref{lem:curvature} bounds its covariance
by $s\Id$. Gaussian differentiation is valid without compact support.
Thus $\mu_s$ is centered and log-concave, with a smooth convex potential.

Set $\beta_s=\sqrt{8s/3}$ and $a_s=a+\beta_s$. For a scalar standard
Gaussian $Z$, $\E\exp(sZ^2/\beta_s^2)=2$. The inequality
\[
 \frac{(u+t)^2}{(a+\beta_s)^2}
 \le\frac{a}{a+\beta_s}\frac{u^2}{a^2}
     +\frac{\beta_s}{a+\beta_s}\frac{t^2}{\beta_s^2}
\]
and H\"older's inequality show that $\mu_s$ satisfies
\eqref{eq:intro-subgaussian} with parameter $a_s$.
Theorem~\ref{thm:main}, with the Sobolev extension in
Section~\ref{sec:calculus}, gives entropy coefficient $Ca_s^2$.
In the conventions of \cite[Definitions~1--2]{OttoVillani2000},
$\Ent(f^2)\le L\int|Df|^2$ is $\mathrm{LSI}(2/L)$, whose
transport conclusion is $W_2^2\le LH$. Their Theorem~1 applies to
the smooth convex potential of $\mu_s$, and hence
\[
 W_2(\nu_s,\mu_s)^2\le Ca_s^2H(\nu_s\mid\mu_s)
 \le Ca_s^2H(\nu\mid\mu).
\]
The last inequality is data processing for the common Gaussian
kernel, as in Section~\ref{sec:calculus}. Coupling a vector with
itself plus independent Gaussian noise gives
$W_2(\mu_s,\mu),W_2(\nu_s,\nu)\le\sqrt{ns}$.
Letting $s\downarrow0$ at fixed $n$, with $a_s\to a$, proves the claim.
\end{proof}
